% ---------
% --------
%!TEX TS-program = pdflatex
%!TEX encoding = UTF-8 Unicode
%!TEX spellcheck = English (Aspell)

\documentclass[11pt]{article}
\usepackage{lucy,handout,graphicx}
\usepackage{fourier-orns}
\usepackage{mathtools}
\usepackage[normalem]{ulem} % either use this (simple) or
\pagestyle{empty}

\allowdisplaybreaks

\begin{document}


\begin{center}
\Large\textbf{CSCI 332, Fall 2026}%
\\
\LARGE\textbf{In-class Activity}%
\end{center}

In this activity, you will analyze the runtime of the gradeschool algorithm
for multiplying two numbers $x$ and $y$, where $x$ has $m$ digits and $y$ has 
$n$ digits.

\begin{enumerate}
   \item 
First, let's consider the grade-school algorithm for multiplication.
This algorithm works by multiplying each digit of the second number
by the entire first number, shifting appropriately, and then summing
the results.

For example, to multiply 5281 by 23492, we would compute (you should fill this
in):

\[
\begin{array}{r}
   \phantom{00000}5281 \\
\times\phantom{0}23492 \\
\hline
\end{array}
\]
\vspace{2in}

\item Check your answer with another group. Do they agree?

\item  What ``milestones'' would be good to count in in this algorithm to get a
sense of its runtime? What is the runtime of this algorithm in terms of $m$ and
$n$? 
\vspace{2in}

\item Is the runtime
different in a best and worst case? Try to give $\Theta$ bounds for each.
\newpage

\item What are $m$ and $n$ in terms of the values of $x$ and $y$? For example,
if $x = 5281$ and $y = 23492$, what are $m=4$ and $n=5$. What is the general formula for $m$ and $n$?

\vspace{1in}

\item In the mid-1950s, Andrei Kolmogorov, one of the giants of 20th century
mathematics, publicly conjectured that there is no algorithm to multiply two
n-digit numbers in subquadratic time (less than $O(n^2)$). Does your group (circle one):
\begin{itemize}
   \item Agree: $O(n^2)$ is the best we can do.
   \item Disagree: There is a better algorithm (for example, $O(n^{1.5})$ or $O(n \log n)$).
   \item We still don't know: no one has been able to come up with a proof either way.
\end{itemize}


\end{enumerate}

%%%%%%%%%%%%%%%%%%%%
\end{document}
